\section*{Bab \ref{ch:b1:replication-mitosis} --- Replikasi DNA dan Mitosis}

\begin{solution}{exo:b1:replication-mitosis:1}
Polimerase hanya menambahkan \omterm{def:b1:nucleic-acids:nucleotide}{nukleotida} pada sebuah ujung $3'$,
sehingga untai baru tumbuh $5' \to 3'$; kedua cetakannya antiparalel,
sehingga pada sebuah garpu satu untai baru dapat tumbuh ke arah
garpunya secara sinambung sedangkan yang lain harus tumbuh menjauhinya,
dengan memulai lagi dalam keping seiring terpaparnya cetakan baru.
\end{solution}

\begin{solution}{exo:b1:replication-mitosis:2}
Helikase, yang mengurai heliksnya; \omterm{def:b1:proteins:peptide}{protein} pengikat untai tunggal, yang
menahan untainya tetap terpisah; topoisomerase, yang melepaskan
pilinan di depannya; primase, yang menaruh primer \omterm{def:b1:nucleic-acids:chain}{RNA}; \omterm{def:b1:replication-mitosis:fork}{DNA polimerase},
yang memperpanjang primernya lalu mengoreksi; sebuah polimerase kedua,
yang mengganti primernya; dan ligase, yang menyegel fragmennya.
\end{solution}

\begin{solution}{exo:b1:replication-mitosis:3}
G$_1$: $q$, 2n \omterm{def:b1:genomes:genome}{kromosom} berkromatid satu. G$_2$: $2q$, 2n \omterm{def:b1:genomes:genome}{kromosom}
berkromatid dua. Sesudah \omterm{def:b1:replication-mitosis:mitosis}{mitosis}: $q$ tiap \omterm{def:b1:cell-unit-of-life:cell}{sel} anak, 2n \omterm{def:b1:genomes:genome}{kromosom}
masing-masing.
\end{solution}

\begin{solution}{exo:b1:replication-mitosis:4}
Profase (pemadatan, \omterm{def:b1:replication-mitosis:mitosis}{gelendongnya} terbentuk), prometafase (selubungnya
rusak, kinetokornya melekat), metafase (penjajaran pada khatulistiwanya),
anafase (\omterm{def:b1:replication-mitosis:mitosis}{kromatidnya} terpisah), telofase (selubungnya terbentuk
kembali), lalu sitokinesis.
\end{solution}

\begin{solution}{exo:b1:replication-mitosis:5}
Separuh \omterm{def:b1:genomes:genome}{genomnya} adalah sintesis untai tertinggal: $\num{6.4e9}/150 =
\num{4.3e7}$ fragmen; sebanyak itu pula primer, penyingkiran dan
penyambungannya.
\end{solution}

\begin{solution}{exo:b1:replication-mitosis:6}
$\num{6.4e9}\times 10^{-5} = \num{64000}$; $\times 10^{-7}$: 640;
$\times 10^{-9}$: sekitar 6 mutasi tiap pembelahan.
\end{solution}

\begin{solution}{exo:b1:replication-mitosis:7}
$5000/80 = 62$ pembelahan. Sebuah tumor harus membelah tanpa batas;
tanpa telomerase \omterm{def:b1:genomes:karyotype}{telomernya} akan habis dan \omterm{def:b1:cell-unit-of-life:cell}{selnya} akan terhenti atau
mati, sehingga hampir semua kanker mengaktifkan kembali \omterm{def:b1:enzymes:enzyme}{enzim} itu.
\end{solution}

\begin{solution}{exo:b1:replication-mitosis:8}
Daurnya $= 1\,\mathrm{h}/0.05 = \qty{20}{h}$; S $= 0.30\times 20 =
\qty{6}{h}$.
\end{solution}

\begin{solution}{exo:b1:replication-mitosis:9}
G$_1$: 8 \omterm{def:b1:genomes:genome}{kromosom}, 8 \omterm{def:b1:replication-mitosis:mitosis}{kromatid}, $q$. G$_2$: 8, 16, $2q$. Metafase:
8, 16, $2q$. \omterm{def:b1:cell-unit-of-life:cell}{Sel} anak: 8, 8, $q$.
\end{solution}

\begin{solution}{exo:b1:replication-mitosis:10}
Putaran baru mulai tiap \qty{20}{min} pada titik awalnya sebelum
putaran sebelumnya berakhir: replikasinya bertindih, dan sebuah \omterm{def:b1:cell-unit-of-life:cell}{sel} yang
baru lahir sudah memiliki \omterm{def:b1:nucleic-acids:chain}{DNA} yang sebagian tereplikasi. Pada
pembelahannya \omterm{def:b1:genomes:genome}{kromosom} yang menyelesaikan putarannya dimulai
\qty{40}{min} lebih dahulu, sebuah putaran kedua yang dimulai
\qty{20}{min} lebih dahulu garpunya berada di tengah, dan putaran ketiga
baru saja mulai: 4 titik awal dan 6 garpu pada \omterm{def:b1:cell-unit-of-life:cell}{sel} yang membelah itu.
\end{solution}

\begin{solution}{exo:b1:replication-mitosis:11}
Tak ada \omterm{def:b1:replication-mitosis:mitosis}{gelendong} yang terbentuk: \omterm{def:b1:genomes:genome}{kromosomnya} memadat tetapi tak dapat
berjajar maupun terpisah, dan \omterm{def:b1:cell-unit-of-life:cell}{selnya} terhenti pada \omterm{prop:b1:replication-mitosis:checkpoints}{titik pemeriksaan}
\omterm{def:b1:replication-mitosis:mitosis}{gelendongnya} dalam keadaan serupa metafase. Bila ia tergelincir kembali
ke \omterm{def:b1:replication-mitosis:cycle}{interfase} tanpa membelah ia berkromosom 4n (tetraploid). \omterm{def:b1:cell-unit-of-life:cell}{Sel} yang
terhenti dengan \omterm{def:b1:genomes:genome}{kromosom} yang memadat dan terpisah justru itulah yang
dibutuhkan sebuah \omterm{def:b1:genomes:karyotype}{kariotipe}: obatnya dipakai untuk menimbunnya.
\end{solution}

\begin{solution}{exo:b1:replication-mitosis:12}
Replikasinya menjamin bahwa tiap \omterm{def:b1:genomes:genome}{kromosom} menjadi dua \omterm{def:b1:replication-mitosis:mitosis}{kromatid} yang
serupa, sampai satu kekeliruan tiap semiliar; \omterm{def:b1:replication-mitosis:mitosis}{gelendongnya} menjamin
bahwa tiap \omterm{def:b1:cell-unit-of-life:cell}{sel} anak menerima satu dari tiap pasangnya, dengan
memautkan tiap \omterm{def:b1:genomes:genome}{kromosom} dari kedua kutubnya dan menahan anafasenya
sampai tiap \omterm{def:b1:genomes:genome}{kromosomnya} tertarik. \omterm{prop:b1:replication-mitosis:checkpoints}{Titik pemeriksaan} G$_2$-nya
memastikan penyalinannya selesai sebelum penghitungannya mulai; \omterm{prop:b1:replication-mitosis:checkpoints}{titik
pemeriksaan} \omterm{def:b1:replication-mitosis:mitosis}{gelendongnya} memastikan penghitungannya siap sebelum
pemisahannya. Sebuah \omterm{def:b1:cell-unit-of-life:cell}{sel} yang salah memilah berakhir dengan satu
\omterm{def:b1:genomes:genome}{kromosom} terlalu banyak atau terlalu sedikit (aneuploidi): sebuah
salinan yang benar sampai hurufnya, yang diantarkan ke alamat yang
keliru --- yang membunuh sebagian besar \omterm{def:b1:cell-unit-of-life:cell}{sel} semacam itu dan menjadi ciri
sebagian besar kanker.
\end{solution}

\begin{solution}{pb:b1:replication-mitosis:1}
\textbf{1.} $50\times 28\,800 = \qty{1.44e6}{bp}$.
\textbf{2.} \qty{2.88e6}{bp}.
\textbf{3.} $\num{6.4e9}/\num{2.88e6} = 2220$ titik awal.
\textbf{4.} \qty{2.9}{Mb}, yakni $2.9\times 10^6\times 0.34\,\mathrm{nm} =
\qty{0.98}{mm}$ \omterm{def:b1:nucleic-acids:chain}{DNA} antara titik awalnya.
\textbf{5.} Jaraknya $\num{6.4e9}/\num{40000} = \qty{160}{kb}$; tiap
garpu menempuh sekitar \qty{80}{kb}, kerja setengah jam.
\textbf{6.} \qty{125}{Mb} tiap garpu pada \qty{50}{bp/s}: $\num{2.5e6}$
s, yakni 29 hari.
\textbf{7.} Kecepatan garpunya dibatasi kimianya dan kromatinnya
(perlambatan dua puluh kali terhadap bakteri); titik awalnya dapat
dilipatgandakan tanpa batas, sehingga fase S-nya dipendekkan dengan
menambah titik awal, bukan dengan memacu garpunya.
\textbf{8.} Sintesis tertinggalnya meliputi separuh \omterm{def:b1:genomes:genome}{genomnya},
$\num{3.2e9}$ \omterm{def:b1:nucleic-acids:nucleotide}{nukleotida}: $\num{2.1e7}$ fragmen (kedua garpu tiap titik
awalnya masing-masing memiliki sebuah untai tertinggal, dan kedua
paruhnya berjumlah senilai satu \omterm{def:b1:genomes:genome}{genom}).
\textbf{9.} $\num{2.1e7} + 2\times\num{40000} = \num{2.1e7}$ primer.
\textbf{10.} Sekitar $\num{2.1e7}$ penyambungan (satu tiap fragmen)
ditambah satu tiap pertemuan garpunya.
\textbf{11.} $2\times\num{6.4e9} = \num{1.28e10}$ \omterm{def:b1:nucleic-acids:nucleotide}{nukleotida}.
\textbf{12.} $\num{2.56e10}$ \omterm{def:b1:water-small-molecules:atp}{ATP}; sepanjang \qty{28800}{s},
$\num{8.9e5}$ tiap sekon.
\textbf{13.} Sekitar \qty{0.1}{\%} perputaran \omterm{def:b1:water-small-molecules:atp}{ATP} \omterm{def:b1:cell-unit-of-life:cell}{selnya}:
pemolimerannya murah; yang mahal adalah membuat \omterm{def:b1:nucleic-acids:nucleotide}{nukleotida} dan
\omterm{def:b1:genomes:nucleosome}{histonnya}.
\textbf{14.} \num{64000}, 640 dan 6.4 kekeliruan.
\textbf{15.} $1500\times 10^{-9} = \num{1.5e-6}$: satu \omterm{def:b1:genomes:gene}{gen} tiap tujuh
ratus ribu tiap pembelahan.
\textbf{16.} $10^{16}\times 6.4 = \num{6.4e16}$ mutasi --- tiap
perubahan tunggal \omterm{def:b1:genomes:genome}{genomnya} yang mungkin, berkali-kali, yang tersebar di
antara $10^{13}$ \omterm{def:b1:cell-unit-of-life:cell}{selnya}: hampir semuanya jatuh pada \omterm{def:b1:nucleic-acids:chain}{DNA} bukan penyandi
atau pada \omterm{def:b1:cell-unit-of-life:cell}{sel} yang ditanggalkan, dan sebuah \omterm{def:b1:cell-unit-of-life:cell}{sel} yang bermutasi adalah
satu di antara miliaran; hanya beberapa gabungan dalam satu \omterm{def:b1:cell-unit-of-life:cell}{sel} yang
penting.
\textbf{17.} Seratus kali ($10^{-7}$): urutan mutasi yang membuat
sebuah kanker, yang memakan waktu berdasawarsa untuk menumpuk pada
$10^{-9}$, menumpuk jauh lebih cepat --- cacat bawaan pada perbaikan
salah pasang menyebabkan kanker usus besar pada awal masa dewasa.
\textbf{18.} $\num{2.3e6}$ pasangan tiap garpu pada \qty{1000}{bp/s}:
\qty{2300}{s}, \qty{38}{min}.
\textbf{19.} Sekitar dua putaran yang bertindih; sebuah \omterm{def:b1:cell-unit-of-life:cell}{sel} yang baru
lahir membawa dua titik awal (yang masing-masing sudah tereplikasi
sekali) dan garpu yang berada di tengah: \omterm{def:b1:genomes:genome}{kromosomnya} sebagian sudah
tereplikasi saat lahir.
\textbf{20.} $\num{4.6e6}/1500 = 3070$ fragmen tiap putaran, sekitar
1.3 tiap sekon.
\textbf{21.} $\num{4.6e6}\times 10^{-9} = \num{4.6e-3}$ kekeliruan tiap
putaran: sekitar satu \omterm{def:b1:cell-unit-of-life:cell}{sel} anak dari 200 membawa sebuah mutasi baru.
\textbf{22.} $10^9\times\num{4.6e-3} = \num{4.6e6}$ mutasi baru tiap
mililiter; sebuah \omterm{def:b1:genomes:gene}{gen} \qty{1}{kb} tertentu terkena $\num{4.6e6}\times
1000/\num{4.6e6} = 1000$ kali.
\textbf{23.} Dengan seribu mutasi bebas pada \omterm{def:b1:genomes:gene}{gen} mana pun tiap
mililiter tiap pembelahan, tiap perubahan tunggal yang mungkin ---
termasuk yang memberikan ketahanan --- sudah hadir dalam sebuah biakan
besar sebelum antibiotiknya dikenakan; obatnya memilih, ia tidak
menciptakan.
\textbf{24.} Manusia: $\num{6.4e9}$ pasangan, \num{40000} titik awal,
\qty{50}{bp/s}, \qty{8}{h}. Bakteri: $\num{4.6e6}$ pasangan, satu titik
awal, \qty{1000}{bp/s}, \qty{40}{min}. Sebuah \omterm{def:b1:genomes:genome}{genom} seribu kali lebih
besar, yang disalin garpu dua puluh kali lebih lambat, hanya dalam dua
belas kali waktunya: selisihnya adalah puluhan ribu titik awal yang
bekerja serentak.
\textbf{25.} Sekitar \num{2200} titik awal paling sedikit, seandainya
semuanya menyala sekaligus; sesungguhnya sekitar \num{40000}, satu tiap
\qty{160}{kb}, yang menyala berurutan sehingga tiap garpunya menempuh
sekitar \qty{80}{kb}.
\end{solution}
